% Generated draft: requires mathematical and source review.
\documentclass[11pt]{article}
\usepackage[margin=25mm]{geometry}
\usepackage{fontspec}
\setmainfont{Latin Modern Roman}
\usepackage{amsmath,amssymb,mathtools,unicode-math}
\setmathfont{Latin Modern Math}
\usepackage{enumitem,needspace,xurl,hyperref}
\setlength{\parindent}{0pt}
\setlength{\parskip}{5pt}
\setlength{\emergencystretch}{3em}
\newcommand{\N}{\mathbb{N}}
\newcommand{\Z}{\mathbb{Z}}
\newcommand{\Q}{\mathbb{Q}}
\newcommand{\R}{\mathbb{R}}
\newcommand{\C}{\mathbb{C}}
\newcommand{\F}{\mathbb{F}}
\newcommand{\A}{\mathbb{A}}
\newcommand{\PP}{\mathbb{P}}
\newcommand{\E}{\mathbb{E}}
\newcommand{\eps}{\varepsilon}
\newcommand{\dd}{\,\mathrm{d}}
\begin{document}
\section*{UnsolvedMath ID: 800688}
\textbf{Source problem code:} catalogue entry
\subsection{Definitions and mathematical statement}
A holomorphic foliation has complex leaves; its leaf dimension is the complex dimension of a leaf on its regular locus. For terminology, in the usual setting of a foliation by complex curves on a complex manifold \(X\) of dimension \(n\), a current \(T\) of bidimension \((1,1)\) acts on compactly supported smooth \((1,1)\)-forms. On the regular locus, such a current is directed by the foliation when \(T\wedge\alpha=0\) for every local holomorphic \((1,0)\)-form \(\alpha\) annihilating the tangent directions of the leaves. It is \(dd^c\)-closed when
\[
 dd^cT=0
\]
in the sense of currents. Here one may use the convention \(dd^c=i\partial\bar\partial\); the vanishing condition is unchanged by a nonzero normalization factor.

\textbf{Problem (Dinh--Sibony, 2018, Problem 15).} For holomorphic foliations whose leaves have complex dimension \(k>1\), identify an appropriate notion of objects replacing the directed \(dd^c\)-closed currents used for foliations by complex curves. These objects should permit the development of an ergodic theory for such foliations.

This is a request for a suitable mathematical framework, rather than a conjecture with a specified conclusion. The supplied statement does not specify the ambient manifold, compactness or metric assumptions, permitted singularities, positivity or finiteness conditions on currents, or the particular ergodic properties the replacement objects should support. The terminology above explains the curve case on the regular locus; it does not supply these missing hypotheses or prescribe a higher-dimensional replacement.
\subsection{Short English statement}
Dinh and Sibony ask what should replace directed \(dd^c\)-closed currents when studying holomorphic foliations with leaves of complex dimension greater than one. The desired objects should support an ergodic theory. The source leaves the precise framework and required ergodic properties unspecified.
\subsection{Sources}
\begin{itemize}
\item Local UnsolvedMath source registry, permanent ID \(800688\), source problem code \texttt{catalogue entry}, title \emph{Holomorphic Foliation Theory --- Problem 15}. The supplied registry labels the problem \texttt{open}; this is a supplied status label, not an independently verified mathematical or literature assessment.
\item Dinh and Sibony, \emph{Some Open Problems on Holomorphic Foliation Theory} (2018), Problem 15. Supplied source URL: \url{https://arxiv.org/abs/1807.05640}. The record attributes the problem to Dinh and Sibony and gives proposed year 2018.
\item Supplied provenance records access on 2026-10-16 and extraction method \texttt{arxiv\_latex}. These are supplied provenance details; the source was not independently accessed or verified for this draft.
\end{itemize}
\end{document}
